\section{Results}\label{sec:results}

\subsection{Template Features}

The Seka template successfully implements all required features for professional academic writing.

\subsection{Figure Integration}

Figure~\ref{fig:workflow} illustrates the document compilation workflow.

\begin{figure}[h]
\centering
\fbox{\parbox{0.8\columnwidth}{
\centering
\includegraphics[width=0.7\columnwidth]{seka/figures/kepler.gl (3).png} % Replace with actual figure
}}
\caption{Document compilation workflow}
\label{fig:workflow}
\end{figure}

\subsection{Table Formatting}

Tables are formatted using the \texttt{booktabs} package for professional quality:

\begin{table}[h]
\centering
\caption{Performance Metrics}
\label{tab:metrics}
\begin{tabular}{lcccc}
\toprule
\textbf{Method} & \textbf{Accuracy} & \textbf{Precision} & \textbf{Recall} & \textbf{F1} \\
\midrule
Method A & 94.2\% & 0.93 & 0.91 & 0.92 \\
Method B & 96.1\% & 0.95 & 0.94 & 0.94 \\
Method C & 95.3\% & 0.94 & 0.92 & 0.93 \\
\bottomrule
\end{tabular}
\end{table}

\subsection{Mathematical Typesetting}

Complex equations are rendered beautifully using \texttt{amsmath}:

\begin{equation}\label{eq:complex}
\frac{\partial \mathbf{u}}{\partial t} + (\mathbf{u} \cdot \nabla) \mathbf{u} = -\frac{1}{\rho} \nabla p + \nu \nabla^2 \mathbf{u}
\end{equation}

\begin{note}
Equation~\ref{eq:complex} shows the Navier-Stokes equations for fluid dynamics, demonstrating support for complex mathematical notation.
\end{note}

\subsection{Code Listings}

Programming code is displayed with syntax highlighting:

\begin{code}
\begin{lstlisting}[language=R, caption={R code for statistical analysis}]
# Load required libraries
library(ggplot2)
library(dplyr)

# Process data
results <- data %>%
    group_by(category) %>%
    summarise(
        mean = mean(value),
        sd = sd(value)
    )
\end{lstlisting}
\label{code:analysis}
\end{code}